\documentclass[a4paper,12pt]{article}
\usepackage[utf8]{inputenc}
\usepackage{calc}
\usepackage{eso-pic}
\usepackage{graphicx}
\usepackage{parskip}
\usepackage{hyperref}
\usepackage[a4paper,left=25 mm,right=25 mm, bottom=25 mm,
    margin=1in,top=22 mm]{geometry} % Adjusting the margins
\usepackage{tikz}
\usepackage{array}
\usepackage{tabularx}
\usepackage{amsmath}
\usepackage{amssymb}
\usetikzlibrary{calc}

\begin{document}

\begin{titlepage}
\begin{tikzpicture}
    [remember picture, overlay]
    \draw[line width = 2pt, black] 
        ($(current page.north west) + (1cm,-1cm)$) 
        rectangle 
        ($(current page.south east) + (-1cm,1cm)$);
\end{tikzpicture}
    \centering
    \vspace*{0 cm}
    \LARGE
    \textbf{Maulana Abul Kalam Azad University of Technology, WB}
    \vspace{0.5cm}
    
    \includegraphics[width=0.2\textwidth]{makaut.png} % Assuming you have the logo image
    \vspace{0.5cm}
    
    \Large
    \textbf{\textcolor{blue!60}{Software Tools and Technology\\
        -: Lab Notebook :-}}
    \vspace{0.5cm}
    
    \large
    \textbf{\textcolor{red}{Group 11}}
    \vspace{1 cm}
    
    \textbf{Repository Link:} \href{https://github.com/ABIR462/group---26/tree/main}{https://github.com/ABIR462/group---26/tree/main}
    \vspace{1cm}
    
    \textbf{\underline{\textcolor{blue!60}{Group Members}}}
    \vspace{0.5cm}

    \normalsize
    \begin{enumerate}
        \item \textbf{Abir Kumar Kayal}\\
              \textbf{Roll No: 30001223034}\\
              \textbf{Department: Bachelor's of Computer application}
        \item \textbf{Amrita Mondal}\\
              \textbf{Roll No: 30059223011}\\
              \textbf{Department: BSc in Forensic Science}
        \item \textbf{Saikat Biswas}\\
              \textbf{Roll No: 30059223028}\\
              \textbf{Department: BSc in Forensic Science}
        \item \textbf{Soumyajit Das}\\
              \textbf{Roll No: 30001223033}\\
              \textbf{Department: Bachelor's of Computer application}
        \item \textbf{Supratim Moulick}\\
              \textbf{Roll No: 30084323024}\\
              \textbf{Department: BSc in IT (Data Science)}
    \end{enumerate}
    \vspace{0.8 cm}
    
    \textbf{Instructor:} \href{mailto:ayan.ghosh@university.edu}{\textcolor{blue}{Ayan Ghosh}}\\
    \vspace{0.2cm}
    
    \textbf{\textit{Date: September 11, 2024}}

\end{titlepage}
\newpage
\begin{tikzpicture}
    [remember picture, overlay]
    \draw[line width = 2pt, black]
        ($(current page.north west) + (1cm,-1cm)$)
        rectangle 
        ($(current page.south east) + (-1cm,1cm)$);
\end{tikzpicture}
\vspace{-2cm}

\centering
\section*{\underline{\Huge\textbf{\textcolor{blue!60}{Index}}}}
\vspace{0.5cm}

\renewcommand{\arraystretch}{2}
\setlength{\tabcolsep}{0pt} 

\begin{tabular}{|>{\centering\arraybackslash}p{80pt}|>{\centering\arraybackslash}p{350pt}|}
\hline
\textbf{Serial No.} & \textbf{Questions} \\
\hline
1 & Introduction to GitHub and GitHub Desktop version installation \\\hline

2 & Building a C program for a calculator in the local repository,committing,and publishing it as a public repository \\\hline
3 & Converting a submit button to Chin Tapak Dum Dum \\\hline
4 & Build a cv in latex\\\hline
5 & Branching and Merging \\\hline


\end{tabular}

\newpage
\begin{tikzpicture}
    [remember picture, overlay]
    \draw[line width = 2pt, black] 
        ($(current page.north west) + (1cm,-1cm)$) 
        rectangle 
        ($(current page.south east) + (-1cm,1cm)$);
\end{tikzpicture}
\vspace{-2cm}
% Lab notebook entries
\section*{\Huge{\textcolor{blue!60}{Lab Notebook Entries}}}

\subsection*{Entry by Abir Kumar Kayal}
\textit{Date: [\today]}\\
\vspace{1 cm}
\begin{figure}[h!]
   \centering
    \includegraphics[width=0.5\linewidth]{GitHub-logo.png}
\end{figure}
\vspace{0.5 cm}

\section*{\Huge{\underline{{-:GitHub:-}}}}
\paragraph{\noindent}{GitHub is a web-based platform that allows developers to host, share, and collaborate on software projects. It provides a version control system powered by Git, enabling teams to track changes, manage code repositories, and work together seamlessly across different locations. GitHub supports collaborative development through features like pull requests, issues, and project boards, making it essential for open-source projects and professional software development. Additionally, it integrates with various development tools, enhancing productivity and streamlining the software development lifecycle.}

\vspace{1 cm}
\subsection*{\underline{-:Installation:-}}
\paragraph{\setlength{\parindent}{0pt}}{Installing GitHub Desktop is a straightforward process that enhances your workflow by providing a user-friendly interface for managing repositories. To begin, download the installer from the [official GitHub Desktop website](https://desktop.github.com/) for your operating system—Windows or macOS. After downloading, run the installer and follow the on-screen instructions to complete the setup. Once installed, launch the application and sign in with your GitHub credentials, or create a new account if needed. GitHub Desktop simplifies the process of cloning repositories, making commits, and managing branches, making it an invaluable tool for developers of all skill levels. For Linux users, alternative methods like using Wine or other Git clients are available.}


\newpage
\begin{tikzpicture}
    [remember picture, overlay]
    \draw[line width = 2pt, black] 
        ($(current page.north west) + (1cm,-1cm)$) 
        rectangle 
        ($(current page.south east) + (-1cm,1cm)$);
\end{tikzpicture}
\vspace{-2cm}
\subsection*{Entry by Amrita mondal}
\textit{Date: [\today]}\\
\subsection{Building a C program for a calculator in the local
repository,committing,and publishing it as a public repository}
\section{Introduction}
This document outlines the process of creating a basic calculator in C, committing it to a local Git repository, and pushing it to a public repository on GitHub.
\section{Step 1: Writing the C Program}
Below is the source code for a simple calculator written in C:
\begin{figure}[h!]
    \centering
    \includegraphics[width=0.8\linewidth]{Screenshot 2024-09-21 181207.png} % Adjusted width
     \hspace{4 cm}
    \caption{Screenshot of the source code.}
\end{figure}
\section{Step 2: Create a New Repository on GitHub}

1. Go to your GitHub profile: \href{https://github.com/ABIR462}{GitHub - ABIR462}.
2. Click on \texttt{New} to create a new repository.
3. Name the repository \texttt{C-Calculator} and ensure it is set to \texttt{Public}.
4. Do not initialize the repository with a \texttt{README}, \texttt{.gitignore}, or license, as these are not needed for the local repository you have created.

\section{Step 3: Push Local Repository to GitHub}

Copy the remote URL of the new GitHub repository. It should look like this:
\begin{lstlisting}[caption={Example remote URL}]
https://github.com/ABIR462/C-Calculator.git
\end{lstlisting}

Add the remote repository to your local Git:
\begin{lstlisting}[language=bash, caption={Adding a remote repository}]
git remote add origin https://github.com/ABIR462/C-Calculator.git
\end{lstlisting}

Push your local commits to the GitHub repository:
\begin{lstlisting}[language=bash, caption={Pushing to GitHub}]
git push -u origin master
\end{lstlisting}

\section{Step 4: Verify on GitHub}

Visit your GitHub profile to ensure that the code is now available in your public repository: \href{https://github.com/ABIR462/C-Calculator}{https://github.com/ABIR462/C-Calculator}.



    
\newpage
\begin{tikzpicture}
    [remember picture, overlay]
    \draw[line width = 2pt, black] 
        ($(current page.north west) + (1cm,-1cm)$) 
        rectangle 
        ($(current page.south east) + (-1cm,1cm)$);
\end{tikzpicture}
\vspace{-2cm}
\subsection*{Entry by Supratim Moulick}
\textit{Date: [\today]}\\
\date{\today}
\floatBarrier 
\begin{center}
\section*{\uline{-:Converting a submit button to "chin tapak dum dum":-}}
\end{center}
\paragraph{}
\vspace{0.5cm}
Steps to Fix the Button and Create a Pull Request\\

1. Clone the Repository:
   - You’ve already cloned the repository using GitHub Desktop, so you should have the project on your local machine.\\
\vspace{1cm}

2. Open the Project:
   - Open the project in your preferred IDE as per the instructions in the readme.md.\\
\vspace{1cm}

3. Locate the Button Code:
   - Search for the code where the submit button is defined. This is typically found in the frontend code of the application. Depending on the technology stack used (e.g., HTML/CSS, React, Angular, etc.), it could be in a file like index.html, App.js, ButtonComponent.js, or a similar file.\\
\vspace{1cm}

4. Rename the Button:
   - You renamed the button to "Chin Tapak Dum Dum". If the button’s size became disproportionate, it’s likely due to styling issues.\\
\vspace{1cm}

5. Test the Application:
   - Run the application again to ensure that the button appears correctly and that there are no additional issues.\\
\vspace{1cm}

6. Commit Your Changes:
   - Once the button looks good, commit your changes. Use descriptive commit messages, for example: Fixed button styling after renaming.

   bash
   git add .
   git commit -m "Fixed button styling after renaming to 'Chin Tapak Dum Dum'"
   \\
\vspace{1cm}
   
7. Push Your Changes:
   - Push your changes to your forked repository on GitHub.

   bash
   git push origin main
   

   (Replace main with the correct branch name if it's different.)\\
\vspace{1cm}
   
8. Create a Pull Request:
   - Go to the original repository on GitHub (the one you cloned from).
   - You should see an option to create a pull request from your forked repository. Follow the instructions to create a pull request with a title and description of what you have done.

   Make sure to mention in the pull request that you have fixed the button styling after renaming it.




\newpage
\begin{tikzpicture}
    [remember picture, overlay]
    \draw[line width = 2pt, black] 
        ($(current page.north west) + (1cm,-1cm)$) 
        rectangle 
        ($(current page.south east) + (-1cm,1cm)$);
\end{tikzpicture}
\vspace{-2cm}
\subsection*{Entry by Saikat biswas}
\textit{Date: [\today]}\\
% Write your lab notebook entry here.



\newpage
\begin{tikzpicture}
    [remember picture, overlay]
    \draw[line width = 2pt, black] 
        ($(current page.north west) + (1cm,-1cm)$) 
        rectangle 
        ($(current page.south east) + (-1cm,1cm)$);
\end{tikzpicture}
\vspace{-2cm}
\subsection*{Entry by Soumyajit Das}
\textit{Date: [\today]}\\
% Write your lab notebook entry here.
\section*{\Huge{Git Assignment 3 : Branching and Merging}}

\paragraph{Objective:} Demonstrate proficiency in Git branching, merging, and conflict resolution.

\vspace{0.5cm}

% Manually resizing the images to fit within the page margins
% Screenshot 1
\begin{figure}[h!]
    \centering
    \includegraphics[width=0.8\linewidth]{123.jpg} % Adjusted width
     \hspace{4 cm}
    \caption{Screenshot of the GitHub repository showing the commit history.}
\end{figure}

% Screenshot 2 
\begin{figure}[h!]
    \centering
    \includegraphics[width=0.8\linewidth]{WhatsApp Image 2024-09-21 at 18.29.56_365e138f.jpg} % Adjusted width
     \hspace{4 cm}
    \caption{Screenshot of the GitHub repository showing the branching.}
\end{figure}
\newpage
\begin{tikzpicture}
    [remember picture, overlay]
    \draw[line width = 2pt, black] 
        ($(current page.north west) + (1cm,-1cm)$) 
        rectangle 
        ($(current page.south east) + (-1cm,1cm)$);
\end{tikzpicture}
\vspace{-1.5cm}
% Screenshot 3

\begin{figure}[h!]
    \centering
    \includegraphics[width=0.8\linewidth]{125.jpg}
       \hspace{4 cm}
    \caption{Repository after deleting branches 'feature-1' and 'feature-2'.}
    \label{fig:enter-label}
\end{figure}
% Screenshot 4
\hspace{6 cm}

\hspace{6 cm}
% Screenshot 5

\newpage
\begin{tikzpicture}
    [remember picture, overlay]
    \draw[line width = 2pt, black] 
        ($(current page.north west) + (1cm,-1cm)$) 
        rectangle 
        ($(current page.south east) + (-1cm,1cm)$);
\end{tikzpicture}
\vspace{0.5 cm}
\subsection*{\Huge{Write-Up: Experience with Git Branching and Merging}}
\hspace{1 cm}
\paragraph {\noindent}{This assignment focused on the use of Git branching and merging to manage features in a collaborative environment. After creating separate branches for feature-1 and feature-2, each was developed independently.
When merging feature-1 into the main branch, there were no conflicts. However, the merge of feature-2 caused a conflict in the shared.txt file. Conflict resolution was done manually, ensuring that the changes from both branches were retained.
The practical experience highlighted the importance of clear commit messages and Git’s branching capabilities for parallel development and conflict management.}


\end{document}
